% ===========================================================================
% Diagramme de cas d'utilisation global — Nova AI
% ===========================================================================
\begin{figure}[H]
\centering
\begin{tikzpicture}[
  font=\footnotesize,
  every node/.style={align=center},
  actor/.style ={font=\footnotesize\bfseries},
  uc/.style    ={ellipse, draw=black!75, fill=cyan!8,
                 minimum width=32mm, minimum height=8mm,
                 font=\scriptsize, inner sep=2pt},
  inc/.style   ={->, thick, draw=black!70, dashed, >=Stealth},
  rel/.style   ={-, thick, draw=black!75},
  sys/.style   ={rectangle, draw=black!60, fill=white,
                 minimum width=88mm, minimum height=110mm,
                 line width=1pt}
]

% System box (Nova AI)
\node[sys] (sys) at (0,0) {};
\node[font=\bfseries, anchor=north] at (sys.north) {\bfseries Système Nova~AI};

% Use-case nodes inside the system, arranged in two columns
% Left column — patient flow
\node[uc] (book)     at (-2.0, 4.2) {Réserver un\\rendez-vous};
\node[uc] (cancel)   at (-2.0, 3.0) {Annuler un\\rendez-vous};
\node[uc] (resched)  at (-2.0, 1.8) {Reprogrammer\\un rendez-vous};
\node[uc] (uploadrx) at (-2.0, 0.6) {Envoyer une\\ordonnance};
\node[uc] (askinfo)  at (-2.0,-0.6) {Demander des\\informations};
\node[uc] (urgent)   at (-2.0,-1.8) {Signaler une\\urgence};

% Right column — staff flow
\node[uc] (viewcal)  at ( 2.0, 4.2) {Consulter le\\calendrier};
\node[uc] (mngsvc)   at ( 2.0, 3.0) {Gérer les\\services};
\node[uc] (mngteam)  at ( 2.0, 1.8) {Gérer l'équipe};
\node[uc] (revenue)  at ( 2.0, 0.6) {Voir les\\indicateurs};
\node[uc] (asknova)  at ( 2.0,-0.6) {Interroger Nova\\(vocal)};
\node[uc] (settings) at ( 2.0,-1.8) {Configurer\\horaires \& IA};

% Shared
\node[uc] (auto-ai)  at ( 0.0,-3.4) {Classifier l'intention\\(IA)};

% Actors outside the system
\node[actor]                                 (patient) at (-7.5, 1.5) {\faUser\ Patient};
\node[actor, font=\footnotesize\bfseries]    (owner)   at ( 7.5, 2.5) {Propriétaire\\(médecin / gérant)};
\node[actor, font=\footnotesize\bfseries]    (recept)  at ( 7.5,-0.5) {Réceptionniste\\/ assistant};
\node[actor, font=\footnotesize\bfseries]    (system)  at ( 0.0,-5.0) {\textit{<<système IA>>}\\Anthropic Claude};

% Patient relationships
\draw[rel] (patient) -- (book);
\draw[rel] (patient) -- (cancel);
\draw[rel] (patient) -- (resched);
\draw[rel] (patient) -- (uploadrx);
\draw[rel] (patient) -- (askinfo);
\draw[rel] (patient) -- (urgent);

% Owner relationships
\draw[rel] (owner) -- (viewcal);
\draw[rel] (owner) -- (mngsvc);
\draw[rel] (owner) -- (mngteam);
\draw[rel] (owner) -- (revenue);
\draw[rel] (owner) -- (asknova);
\draw[rel] (owner) -- (settings);

% Receptionist subset
\draw[rel] (recept) -- (book);
\draw[rel] (recept) -- (resched);
\draw[rel] (recept) -- (viewcal);

% AI inclusion
\draw[inc] (book)    -- node[midway, sloped, above=-1pt, font=\tiny]{<<include>>} (auto-ai);
\draw[inc] (askinfo) -- node[midway, sloped, above=-1pt, font=\tiny]{<<include>>} (auto-ai);
\draw[inc] (urgent)  -- node[midway, sloped, above=-1pt, font=\tiny]{<<include>>} (auto-ai);
\draw[inc] (asknova) -- node[midway, sloped, above=-1pt, font=\tiny]{<<include>>} (auto-ai);
\draw[rel] (system)  -- (auto-ai);

\end{tikzpicture}
\caption{Diagramme de cas d'utilisation global : trois acteurs humains (patient, propriétaire, réceptionniste) interagissent avec le système, qui s'appuie sur l'acteur secondaire \textit{<<système IA>>} Anthropic Claude pour classifier les intentions et conduire les conversations multi-tours.}
\label{fig:uc}
\end{figure}
